\makeatletter
\providecommand{\setprimedlabel}[1]{\protected@edef\@currentlabel{#1}}
\makeatother
\section{Five estimates at a least counterexample}
\label{sec:comparison-lemmas}

\paragraph{What this section proves and how to read it.}
The comparison theorem (\cref{thm:comparison}) reduces the baseline theorem to the
inequality $\Theta<1$ between finitely many prices. Its proof, in the next
section, composes five inequalities proved here at a least normalized
counterexample: the first comparison $Y_\lambda\ge(1-L_0-o(1))Y$
(\cref{lem:record:first-stage}); the reserved cost
$|Y_\nu-Y_\lambda|\le(\mathcal E+o(1))Y$, with $\mathcal E$ a complete
polynomial in the prices and the auxiliary prices
(\cref{lem:record:record-expansion}); the wrap
$|J_\nu-2^kY_\nu|\le(2^kW+o(1))Y$ (\cref{lem:record:wrap-expansion}); the
index-support inequality
$J_\nu\le Y_\nu+\frac12[(V_{\rm L}-e_{\rm L})+(V_{\rm R}-e_{\rm R})]$
(\cref{lem:record:common-mean,cor:index-support-frame}); and the envelope
bound $V_{\rm L},V_{\rm R}\le(2^kJ+o(1))Y$ (\cref{lem:envelope-energy}). They
rest on one calculus of energies under the reserved Gram bracket
(\cref{lem:sharp-first,lem:bracket-energy,lem:energy-toolkit}), which Part~V
uses as well: for the variance credits and envelope class bounds of (CP5),
together with the variance form of \cref{lem:record:variance-interface}
(\cref{lem:var-gain,prop:background-variance,lem:class-credit,thm:envelopes,cor:envelope-costs}),
for the mean credit (CP7) (\cref{prop:record:mean-credit}), and in the
closing, which evaluates the complete polynomial (\cref{cert:closing}). Every
statement uses only the frame data and the hypotheses of
\cref{thm:comparison} that it lists, and Part~II; none refers to a particular
local law.

\paragraph{How the inequalities fit together.}
The reference $\mu$ is reflection positive at every cut and defines
$Y=Y_\mu$. The first-stage law $\lambda$ is a law that dominates $\mu$ in the
cut order at one cut of the family, plus record differences grouped in tags,
each sharp at the cut of its tag, so it loses only $L_0$. The comparison law
$\nu$ accepts every query and is a finite signed combination of nonnegative
accepting laws, but it is not reflection positive; the reserved Gram bracket
controls it record by record. The signed index form $J_\nu$ of
\eqref{eq:record:expansion} carries the constant coefficient $2^k$ on
$Y_\nu$, and every other term of its expansion couples a nonempty set of
queries and cancels by the actual-square kernel unless all its intervals
wrap. On the other side, $J_\nu$ is bounded by the random-inclusion energies
of the matching envelopes. Decompose by the exact support of the index
functions: the component of $J_\nu$ with empty support is $Y_\nu$ itself,
and that of each energy is the mean energy of the envelope, which dominates
$Y_\lambda$. Keeping these two components exactly, rather than bounding
$J_\nu$ by the full envelope energies without subtracting their mean
energies, saves one of the $2^k$ copies of the reserved
comparison cost: in $\Theta$ the coefficient of $\mathcal E$ is $1-2^{-k}$,
while the matching term and the wrap keep coefficient one
(\cref{lem:record:domination}). The index step bounds each half-function
only under its own envelope and uses no symmetry of $\nu$ and no odd-skew
cancellation; skewness enters only through the exact identity of the empty
component.

\paragraph{Conventions.}
Notation, energies $Y_\kappa$, the least normalized counterexample and $o(1)$
are those of the previous section: atoms are written $\Blk$,
$B=e^{U_0}+R_{\rm err}$, $\Nbr_\Blk=\Nbr_{0,\Blk}=\lambda_\Blk=\mu_\Blk$ off the
first-stage atoms, $d=|V|=72$, $(\Lsc,A)$ is a least normalized
counterexample with $\Lsc$ larger than the thresholds of the
lemmas used in this section (every $o(1)$ is as $\Lsc\to\infty$)
and $Y=Y_\mu>0$. A positive measure on the
labels $X_X^V$ of a finite set $X$ of atoms is \emph{admissible} in the sense
of \cref{def:admissible}: its one-slot marginals are multiples of the uniform
law, it does not depend on the indices, and every sharp bound applied to it
is uniform in the shared labels (on one atom, its normalization is a positive
record). Prices are those of \cref{def:price},
and $C_X=\prod_{\Blk\in X}q_\Blk^{2a}$ is the price of selecting every atom of
$X$.

\paragraph{Remainders.}
Every bound below is robust in the sense of \cref{def:robust}: its $o(1)$ is
a finite sum, with coefficients bounded independently of $\Lsc$ and $A$, of
\emph{energy errors} (the relative error of
\cref{prop:record:density-structured}, multiplying a structured-energy bound
once), \emph{cancellation remainders} (integrals of the form bounded in
\cref{lem:record:density-can}) and \emph{large-product remainders} (the terms
of an expansion whose union $E$ of selected atoms has $q_E>N^\tau$, bounded
together by \cref{lem:record:density-lp} with the post-reduction prices of
the selected laws as coefficients). Each statement lists which kinds occur;
this is what \cref{lem:record:rec}(ii) uses.

\paragraph{Fixed selections and summed remainders.}
An individual asymptotic bound below with a selected set $X$ is a bound
for fixed $X$; its large-product remainder need not be uniform as $X$
varies. In a sum over selections we first expand the outside laws and
split by the \emph{full} union $E$ of selected atoms. On $q_E\le N^\tau$
the structured-energy error is uniform in $E$ and its labels. All terms
with $q_E>N^\tau$ are instead summed by
\cref{lem:record:density-lp}, with their post-reduction coefficients.
Thus no summed argument below uses a pointwise $o(1)$ for an arbitrary
large selected set. For fixed $X$, the same procedure also proves the
individual asymptotic statements: if $Z$ is the further selected set,
then
\[
 \sum_{Z:\,q_Xq_Z>N^\tau}\prod_{\Blk\in Z}\eta_\Blk
 \le N^{-\tau/4}q_X^{1/4}
 \prod_\Blk(1+\eta_\Blk q_\Blk^{1/4}).
\]
The factor $q_X^{1/4}$ is fixed in that assertion. When $X$ varies in a
priced expansion it is included in the full union before Rankin's
inequality is applied, and (CP8)(c) bounds the resulting product.

\subsection{Bounding a bilinear form by energies of random restrictions}

The small primes enter the comparison only through the external index. The
next lemma is the inequality behind every index estimate: the one-cut
coupling inequality, the bound of the signed index form by random-inclusion
energies, and the common mean are its cases.

\begin{lemma}[Index-support lemma]
\label{lem:record:common-mean}
Let $\mathcal X_+$ be a finite set, and $M_{\rm L}$, $M_{\rm R}$, $\Gamma$ real
matrices indexed by $\mathcal X_+\times\mathcal X_+$, with $M_{\rm L},M_{\rm R}$
symmetric and
\begin{equation}\label{eq:record:joint-gram}
 \begin{pmatrix}M_{\rm L}&\Gamma\\ \Gamma^{\mathsf T}&M_{\rm R}\end{pmatrix}
 \succeq0 .
\end{equation}
Let $\Omega=\prod_{j\in\mathcal K}\Omega_j$ be a finite product of finite sets
with the uniform probability measure. On coordinate $j$, let
$E_{j,0}\le\cdots\le E_{j,h_j}=I$ be a chain of conditional-expectation
projections, with $h_j\ge1$ and $E_{j,0}$ the projection onto constants. Put
$\Pi_{j,0}=E_{j,0}$, $\Pi_{j,l}=E_{j,l}-E_{j,l-1}$ for $l\ge1$, and
\[
 H_{\boldsymbol l}=\bigotimes_{j\in\mathcal K}\operatorname{ran}\Pi_{j,l_j},
 \qquad 0\le l_j\le h_j.
\]
Choose tails $1=\delta_{j,0}\ge\delta_{j,1}\ge\cdots\ge\delta_{j,h_j}\ge0$
and put $\delta_{\boldsymbol l}=\prod_j\delta_{j,l_j}$. Let $\mathcal T$ on
$L^2(\Omega)$ preserve every $H_{\boldsymbol l}$, act as the identity on
$H_{\boldsymbol0}$, and have norm at most $\delta_{\boldsymbol l}$ on
$H_{\boldsymbol l}$. It acts componentwise on vector-valued maps.
Independently of $\omega$, reveal independent depths $J_j\in\{0,\ldots,h_j\}$
with $\Pr(J_j\ge l)=\delta_{j,l}$, write
$\Jc=(J_j)_j$, and set $Z^\Jc=(\bigotimes_jE_{j,J_j})Z$.
For $X,Y\colon\Omega\to\R^{\mathcal X_+}$ put
\begin{gather*}
 \mathsf J=\E_\omega\ip{X(\omega)}{\Gamma(\mathcal TY)(\omega)},\qquad
 e_{\rm L}=\ip{\E X}{M_{\rm L}\,\E X},\qquad
 e_{\rm R}=\ip{\E Y}{M_{\rm R}\,\E Y},\\
 V_{\rm L}=\E_\Jc\E_\omega\ip{X^\Jc(\omega)}{M_{\rm L}X^\Jc(\omega)},\qquad
 V_{\rm R}=\E_\Jc\E_\omega\ip{Y^\Jc(\omega)}{M_{\rm R}Y^\Jc(\omega)} .
\end{gather*}
Then $V_{\rm L}\ge e_{\rm L}\ge0$, $V_{\rm R}\ge e_{\rm R}\ge0$, and
\begin{enumerate}[label=(\alph*)]
\item
 \begin{equation}\label{eq:record:common-mean}
  \bigl|\mathsf J-\ip{\E X}{\Gamma\,\E Y}\bigr|
  \le\tfrac12\bigl[(V_{\rm L}-e_{\rm L})+(V_{\rm R}-e_{\rm R})\bigr];
 \end{equation}
\item $|\mathsf J|\le V_{\rm L}^{1/2}V_{\rm R}^{1/2}$.
\end{enumerate}

The whole-coordinate case is $h_j=1$: identify $\Jc$ with the set of
coordinates whose depth is one. Then $Z^\Jc=\E[Z\mid\omega_\Jc]$, each
coordinate is included independently with probability $\delta_j=\delta_{j,1}$,
and the sector indexed by $S=\{j:l_j=1\}$ is
$H_S=\bigotimes_{j\in S}L^2_0(\Omega_j)\otimes\bigotimes_{j\notin S}\R\one$,
with weight $\delta_S=\prod_{j\in S}\delta_j$.
\end{lemma}

\begin{proof}
\emph{Factorization.} The diagonal blocks of \eqref{eq:record:joint-gram}
are positive semidefinite, and $\Gamma=M_{\rm L}^{1/2}CM_{\rm R}^{1/2}$ with
$\|C\|\le1$ (\cref{lem:record:block-gram}). Put
$\|Z\|_M^2=\E_\omega\ip{Z(\omega)}{MZ(\omega)}$ for symmetric $M\succeq0$.
Matrices acting on $\R^{\mathcal X_+}$ commute with $\mathcal T$ and with
(conditional) expectations over $\omega$.

\emph{Decomposition.} Nested orthogonal projections satisfy
$E_{j,l}E_{j,l'}=E_{j,\min(l,l')}$. Their differences $\Pi_{j,l}$ are
therefore mutually orthogonal projections summing to $I$. Tensorizing gives
$L^2(\Omega)=\bigoplus_{\boldsymbol l}H_{\boldsymbol l}$, so componentwise
$Z=\sum_{\boldsymbol l}Z_{\boldsymbol l}$ and
$Z_{\boldsymbol0}=\E Z$. These components are orthogonal for every form
$\E_\omega\ip{\cdot}{M\,\cdot}$ with $M\succeq0$, and
\[
 Z^\Jc=\sum_{\boldsymbol l:\,l_j\le J_j\ \text{for every }j}Z_{\boldsymbol l}.
\]
This is the depth version of the exact-support decomposition of
\cref{lem:sector}.

\emph{The form and its energy envelope.} Since $\mathcal T$ preserves the
sectors, only equal sectors pair in $\mathsf J$. On the sector
$H_{\boldsymbol l}$ the factorization gives
\[
 \bigl|\E_\omega\ip{X_{\boldsymbol l}}{\Gamma\mathcal TY_{\boldsymbol l}}\bigr|
 =\bigl|\E_\omega\ip{M_{\rm L}^{1/2}X_{\boldsymbol l}}
 {(C\otimes\mathcal T)M_{\rm R}^{1/2}Y_{\boldsymbol l}}\bigr|
 \le\delta_{\boldsymbol l}\|X_{\boldsymbol l}\|_{M_{\rm L}}
 \|Y_{\boldsymbol l}\|_{M_{\rm R}}.
\]
The zero-sector term is $\ip{\E X}{\Gamma\,\E Y}$. The same coefficient
$\delta_{\boldsymbol l}=\Pr(J_j\ge l_j\ \text{for every }j)$ is the
probability that the random projection retains this sector. Orthogonality
therefore gives the energy identity
\[
 V_{\rm L}=\sum_{\boldsymbol l}\delta_{\boldsymbol l}
 \|X_{\boldsymbol l}\|^2_{M_{\rm L}}
 =e_{\rm L}+\sum_{\boldsymbol l\ne\boldsymbol0}\delta_{\boldsymbol l}
 \|X_{\boldsymbol l}\|^2_{M_{\rm L}},
\]
and likewise for $V_{\rm R}$. In the whole-coordinate case this is
\cref{lem:sector}(d).

\emph{Conclusion.} For (a), sum the bound for the form over
$\boldsymbol l\ne\boldsymbol0$ and use $2xy\le x^2+y^2$. For (b), sum over all
$\boldsymbol l$ and apply Cauchy--Schwarz with weights
$\delta_{\boldsymbol l}$. The energy identity also gives
$V_s\ge e_s\ge0$ for $s\in\{{\rm L},{\rm R}\}$.
\end{proof}

\begin{remark}[The inclusion and depth laws]
In the whole-coordinate special case, the inclusion law makes the energy
exactly the norm
envelope of the index operator: $\Pr(S\subseteq\Jc)=\delta_S$ turns the bound
$\delta_S$ on $H_S$ into the weight of $H_S$ in $V_{\rm L}$ and $V_{\rm R}$.
The proof uses only $\Pr(S\subseteq\Jc)\ge\delta_S$, and among such laws the
independent one minimizes every moment later paid for the inclusions: for
$w_j\ge1$, $\E\prod_{j\in\Jc}w_j=\sum_S\Pr(S\subseteq\Jc)\prod_{j\in S}(w_j-1)
\ge\prod_j(1+\delta_j(w_j-1))$.
For the finite-depth law of (CP0)(d), independence instead gives
$\Pr(J_j\ge k_j\ \text{for every }j)=\prod_j\delta_{j,k_j}$.
These joint tails give the sector weights in the filtration form of the
lemma and the conductor moment $R_I$ of \cref{lem:hyb:moment}.
\end{remark}

The instructive base case is the plain coupling inequality, in which one
positive form plays all three roles.

\begin{corollary}[One-cut coupling inequality]\label{lem:hyb:coupling}
In the whole-coordinate case of \cref{lem:record:common-mean}, let
$C\succeq0$ be symmetric,
$\|F\|_C^2=\ip F{CF}$, and let $T_j$ be doubly stochastic matrices on
$\Omega_j$ with $\|T_jf\|\le\delta_j\|f\|$ for $f$ of mean zero. Let $\omega$
be uniform, $\Pr(\omega'=\eta\mid\omega)=\prod_jT_j(\omega_j,\eta_j)$, and $\Jc$
independent of $(\omega,\omega')$. With
$\mathsf I=\E\ip{X(\omega)}{CY(\omega')}$ and
$\mathsf V=\frac12\E_\Jc\E_\omega[\|X^\Jc\|_C^2+\|Y^\Jc\|_C^2]$,
\begin{equation}\label{eq:hyb:coupling}
 \tfrac12\|\E X+\E Y\|_C^2-\mathsf V\ \le\ \mathsf I\ \le\
 \mathsf V-\tfrac12\|\E X-\E Y\|_C^2 ;
\end{equation}
in particular $|\mathsf I|\le\mathsf V$. No relation between $\E X$ and $\E Y$
is assumed.
\end{corollary}

\begin{proof}
$\E[Y(\omega')\mid\omega]=(TY)(\omega)$ with $T=\bigotimes_jT_j$. Each $T_j$
fixes the constants (its rows sum to one) and maps $L^2_0(\Omega_j)$ into
itself (its columns sum to one), so $T$ satisfies the hypotheses on
$\mathcal T$ (\cref{lem:sector}(a)). Apply (a) with
$M_{\rm L}=M_{\rm R}=\Gamma=C$, for which \eqref{eq:record:joint-gram} holds as
$\bigl(\begin{smallmatrix}C&C\\C&C\end{smallmatrix}\bigr)=
\bigl(\begin{smallmatrix}1&1\\1&1\end{smallmatrix}\bigr)\otimes C$: since
$\mathsf V=\frac12(V_{\rm L}+V_{\rm R})$, it gives
$|\mathsf I-\ip{\E X}{C\,\E Y}|\le\mathsf V-\frac12(\|\E X\|_C^2+\|\E Y\|_C^2)$,
which is \eqref{eq:hyb:coupling}.
\end{proof}

\begin{corollary}[The signed index form and the common mean]
\label{cor:index-support-frame}
Assume (CP0)(c),(d) of \cref{thm:comparison}. Let
$\nu=\bigotimes_\Blk\nu_\Blk$ be a product of finite signed measures on the
$X_\Blk^V$ (for instance the comparison law of (CP4)), and let
$M_{{\rm L},\Blk}$, $M_{{\rm R},\Blk}$ be nonnegative measures with
$\bigl(\begin{smallmatrix}M_{{\rm L},\Blk}&\nu_\Blk\\ \nu_\Blk^*&M_{{\rm R},\Blk}
\end{smallmatrix}\bigr)\succeq_\iota0$ for every atom, as in (CP6). Let $J_\nu$
be the signed index form of \eqref{eq:record:expansion}, $X_\omega$,
$Y_\omega$, $\bar X$ the half-functions of \cref{def:hyb:coupled},
$e_s=Q_{\otimes M_s}(\bar X,\bar X)$ for $s\in\{{\rm L},{\rm R}\}$, and
$V_{\rm L}=V_{\otimes M_{\rm L}}(X)$, $V_{\rm R}=V_{\otimes M_{\rm R}}(Y)$ the
single-half random-inclusion energies of \cref{def:hyb:coupled}. Then:
\begin{enumerate}[label=(\roman*)]
\item $\mathcal T=\bigotimes_e(2P_e-T_e)$, and likewise $\mathcal T^*$, maps
 every exact-conductor sector $H_S$ of the index space into itself, with
 $S$ a vector of conductor depths, is the identity on
 $H_{\boldsymbol0}$ and has norm at most $\delta_S$ on $H_S$;
\item $\E_\omega X_\omega=\E_\omega Y_\omega=\bar X$, and for every value of
 $\Jc$, $X^\Jc_\omega$ carries $g^{I_{b_e},D_{\Jc_e},s_e\bmod D_{\Jc_e}}$ at $u_e$
 and $F_A$ at the other slots of $V_+$, likewise $Y^\Jc_\omega$ at the slots
 $\iota v_e$;
\item $V_s\ge e_s\ge0$,
 \begin{equation}\label{eq:index-support-frame}
  |J_\nu-Y_\nu|\le\tfrac12\bigl[(V_{\rm L}-e_{\rm L})+(V_{\rm R}-e_{\rm R})\bigr],
 \end{equation}
 and $|J_\nu|\le V_{\rm L}^{1/2}V_{\rm R}^{1/2}$.
\end{enumerate}
\end{corollary}

\begin{proof}
(i) The coordinates are the interval and the full prime-power
residues of \cref{def:hyb:coupled}. Convolution preserves each exact
conductor space, as does its adjoint. The interval has the two-sector
chain of constants and their orthogonal complement, with tail one.
On a conductor vector $S$, put $S_e$ for its coordinates belonging to
query $e$. If $S_e=\boldsymbol0$, then $P_e=T_e=I$ on this sector.
Otherwise $P_e=0$ and $2P_e-T_e=-T_e$, whose norm is bounded by the
product of the tails belonging to $S_e$. Tensoring over queries proves
(i), for both $\mathcal T$ and its adjoint. For whole-coordinate
Bernoulli inclusion the same argument is the ordinary support proof.

(ii) For uniform $(b,s)$, $\E_{b,s}g^{I_b,\Mzero,s}=F_A$, and conditioning a
uniform residue on its revealed depths $\Jc_e$ turns $g^{I,\Mzero,s}$ into
$g^{I,D_{\Jc_e},s\bmod D_{\Jc_e}}$, with $D_{\Jc_e}$ the product of the
revealed prime powers
(\cref{lem:twointerval}). The queries have pairwise disjoint endpoints and
independent indices, so $X_\omega$ and $Y_\omega$ are products of functions of
the independent $\omega_e$ at distinct slots, with $F_A$ at the remaining
slots; expectations and conditional expectations factorize over $e$, and the
interval coordinates lie in $\Jc$ with probability one.

(iii) Let $M_{\rm L}$, $M_{\rm R}$, $\Gamma$ be the cut matrices at $\iota$ of
$\bigotimes_\Blk M_{{\rm L},\Blk}$, $\bigotimes_\Blk M_{{\rm R},\Blk}$ and $\nu$;
by \eqref{eq:cc:kron} and the tensorization of brackets
(\cref{lem:record:block-gram}) they satisfy \eqref{eq:record:joint-gram}. Then
$\mathsf J=J_\nu$, $\ip{\bar X}{\Gamma\bar X}=Y_\nu$, and the depth envelope of
(CP0)(d) is that of the filtration form of \cref{lem:record:common-mean}; apply that lemma with (i)
and (ii).
\end{proof}

\begin{remark}[Where the gain comes from]
With $M_{\rm L}=M_{\rm R}=\Gamma\succeq0$, \cref{lem:record:common-mean}(a) is
the plain inequality $|\mathsf J-e|\le V-e$ for the common mean energy $e$ of
one positive form, and gives nothing new: the empty components on the two
sides coincide. The gain comes from envelopes whose mean energies
$e_{\rm L},e_{\rm R}$ exceed the mean energy $Y_\nu$ of the law itself. Only
single-half energies occur, $X$ under $M_{\rm L}$ and $Y$ under $M_{\rm R}$,
which is the order in which \eqref{eq:record:joint-gram} holds; the bracket
with the envelopes exchanged is neither proved nor needed.
\end{remark}

\subsection{The signed average over intervals: cancellation except when intervals wrap}

Recall from \cref{def:hyb:coupled} that $P_e$ averages over the
whole index of the query $e$ and that $T_e$ is the coupled
transition to the next interval with a square residue increment.
The signed index form uses the signed combination $2P_e-T_e$ in
place of $T_e$. Expanding the product over the $k$ queries
\eqref{eq:record:expansion} isolates the term $2^kY_\nu$, in which
every query is fully averaged. Every other term applies the actual
coupled transition $T_e$ to a nonempty set of queries; such a term
pairs two points in ordered intervals forced into a square class,
so it is negligible by the kernel estimate---unless an interval
wraps ($b_e=\Kint$, \cref{subsec:cp:index}), where the two intervals
are no longer ordered and cancellation fails. The lemma bounds the
wrapping terms by the quantity $W$. The signed combination is
controlled by the same $2\times2$ domination as the ordinary
coupling (\cref{cor:index-support-frame}), through
\cref{lem:envelope-energy} below, whose proof does not depend on
this lemma.

\begin{lemma}[Expansion and wrap]\label{lem:record:wrap-expansion}
Assume (CP0)--(CP4) and (CP8) of \cref{thm:comparison}, and put
\[
 W=Be^{S+A}\Bigl[\Bigl(1+\frac{R_I}{2\Kint^{1-2\theta/d}}\Bigr)^k-1\Bigr].
\]
Then $|J_\nu-2^kY_\nu|\le(2^kW+o(1))Y$. The $o(1)$ consists of at most
$3^kn_\nu$ cancellation remainders, with coefficients at most
$2^k\mathfrak v_\nu e^{P+S}$, and of the energy errors and large-product
remainders of the wrapping terms.
\end{lemma}

\begin{proof}
By \eqref{eq:record:expansion},
$J_\nu-2^kY_\nu=\sum_{\mathcal S\ne\varnothing}(-1)^{|\mathcal S|}
2^{k-|\mathcal S|}I_{\nu,\mathcal S}$, where $I_{\nu,\mathcal S}$ carries the
ordinary coupling on the queries of $\mathcal S\subseteq\mathcal E_k$ and the
full average on the others. Split
$I_{\nu,\mathcal S}=\sum_{\mathcal W\subseteq\mathcal S}I_{\nu,\mathcal S,\mathcal W}$
by the set $\mathcal W=\{e\in\mathcal S:b_e=\Kint\}$ of wrapping active
queries; there are $3^k$ pairs $\mathcal W\subseteq\mathcal S$.

\emph{The mass of $\nu$.} On a first-stage atom
$\lambda_\Blk(1)=1+(\rho_\Blk-\mu_\Blk)(1)\le1+p_\Blk$, as the record
differences of (CP2) have mass zero. The cut matrix of
$\zeta^{\rm a}_{\Blk,t}$ is skew, so $\zeta^{\rm a}_{\Blk,t}(1)=0$, and the
bracket of $\zeta^{\rm s}_{\Blk,t}$ in (CP4) at the all-ones vector gives
$|\zeta^{\rm s}_{\Blk,t}(1)|\le\tilde\pi^{\rm s}_{\Blk,t}$. As $q_\Blk^{2a}\ge1$,
$0<\nu(1)\le\prod_\Blk(1+p_\Blk+s_\Blk)\le e^{P+S}$. Hence, every factor
$|\nu|_\Blk/\nu_\Blk(1)$ of (CP4) being at least one,
$\prod_{\Blk\in\mathcal A_N}|\nu|_\Blk\le\mathfrak v_\nu\,\nu(1)\le
\mathfrak v_\nu e^{P+S}$.

\emph{Nonwrapping terms.} If $\mathcal W\ne\mathcal S$, pick
$e\in\mathcal S\setminus\mathcal W$. Then $b_e\le\Kint-1$, the inputs
$g^{I_{b_e},\Mzero,s_e}$ at $u_e$ and $g^{I_{b_e+1},\Mzero,s_e+r_e}$ at $v_e$
lie on adjacent ordered intervals, and $r_e$ is a nonvanishing square
class modulo $\Mzero$, by (CP0)(d) and \cref{lem:imp:padic-envelope}. Expand every
$\nu_\Blk$ into its components: $\nu$ becomes a sum of at most $n_\nu$
product components $\bigotimes_\Blk c_{\Blk,i_\Blk}\nu_{\Blk,i_\Blk}$. By
(CP4)(a) the normalized factors $\bar\nu_{\Blk,i_\Blk}$ of each product
component accept every query, have uniform one-slot marginals and satisfy
(R4) and (R5) of \cref{hyp:record:reference}; their product is a probability
measure of the form required in \cref{lem:record:density-can}, and these are
the hypotheses of that lemma for $(u_e,v_e)$; the other slots carry inputs of
\cref{def:record:functional}(f). At fixed indices each product component
contributes $\prod_\Blk c_{\Blk,i_\Blk}\nu_{\Blk,i_\Blk}(1)$ times an integral
bounded in that lemma, uniformly in the indices, so the bound persists after
integration against their conditional law. Summing the absolute values of the
coefficients over the product components gives
$\prod_{\Blk\in\mathcal A_N}|\nu|_\Blk$. So $I_{\nu,\mathcal S,\mathcal W}$,
with its expansion coefficient, is a sum of at most $n_\nu$ cancellation
remainders whose coefficients add up to at most
$2^k\prod_\Blk|\nu|_\Blk\le2^k\mathfrak v_\nu e^{P+S}$. No other step of this
proof uses the components: the wrapping terms below are bounded through the
brackets of (CP3) and (CP4), which hold for the signed $\nu_\Blk$, and not
through the variation of $\nu$.

\emph{Wrapping terms.} Let $\mathcal W=\mathcal S$ and $j=|\mathcal S|$. All
$j$ intervals wrap with probability $\Kint^{-j}$. Given this, the averages
$P_e$ ($e\notin\mathcal S$) replace those inputs by $F_A$ on both sides, and
the term is $\Kint^{-j}\E_s\ip{X'(s)}{\Gamma(T'Y')(s)}$: $X'(s)$ carries
$g^{I_\Kint,\Mzero,s_e}$ at $u_e$ and $Y'(s)$ carries $g^{I_1,\Mzero,s_e}$ at
$\iota(v_e)$ ($e\in\mathcal S$), both carry $F_A$ elsewhere, $s$ is uniform,
$\Gamma=C_\iota(\nu)$, and $T'$ is the residue part of the coupling on
$\mathcal S$, a tensor product of doubly stochastic kernels with the norms
given by the conductor tails of (CP0)(d) on exact-conductor spaces. On every atom the brackets of
(CP3) and (CP4) add up to
\[
 \begin{pmatrix}M_\Blk&\nu_\Blk\\ \nu_\Blk^*&M_\Blk\end{pmatrix}
 \succeq_\iota0,\qquad
 M_\Blk=\Nbr_\Blk+\sum_t\bigl(\tilde\pi^{\rm s}_{\Blk,t}D^{\rm s}_{\Blk,t}
 +\tilde\pi^{\rm a}_{\Blk,t}D^{\rm a}_{\Blk,t}\bigr),
\]
and these tensorize (\cref{lem:record:block-gram}). As in the proof of
\cref{lem:hyb:coupling}, $T'$ satisfies the hypotheses of
\cref{lem:record:common-mean}, and part (b) bounds the term by
$\Kint^{-j}(V'_{\rm L}V'_{\rm R})^{1/2}$, with $V'_{\rm L},V'_{\rm R}$ the
single-half random-inclusion energies of $X'$ and $Y'$ under
$\bigotimes_\Blk M_\Blk$. Here $M_\Blk-\Nbr_\Blk$ is a nonnegative combination of
the admissible positive reserved-diagonal laws $D^\epsilon_{\Blk,t}$, of total
price $S+A$, so \cref{lem:envelope-energy}, with $T_Q=S+A$ and the fixed
intervals $I_\Kint$, $I_1$ ($\kappa_e=\Kint$, \cref{rem:wrap-coarse}), gives
$V'_{\rm L},V'_{\rm R}\le(Be^{S+A}(\Kint^{2\theta/d}R_I)^j+o(1))Y$. Hence
$|I_{\nu,\mathcal S,\mathcal S}|\le(Be^{S+A}(R_I\Kint^{2\theta/d-1})^j+o(1))Y$.

\emph{Summation.} Summing $2^{k-j}$ times this over the $\binom kj$ sets of
size $j$,
\[
 \sum_{j=1}^k\binom kj2^{k-j}Be^{S+A}\bigl(R_I\Kint^{2\theta/d-1}\bigr)^j
 =Be^{S+A}\bigl[(2+R_I\Kint^{2\theta/d-1})^k-2^k\bigr]=2^kW .\qedhere
\]
\end{proof}

\begin{remark}[The wrap uses the coarse interval factor]
\label{rem:wrap-coarse}
In the wrapping terms the intervals are fixed, so each changed index slot
has the coarse factor $\kappa_v=\Kint$ of \cref{def:record:structured-energy}
($\alpha_I\le\Kint\alpha$), which \cref{prop:record:density-structured} charges
as $\Kint^{\theta/d}$. An active query changes two slots of each single-half
energy, so the wrap pays $\Kint^{2\theta/d}$ per active query; at $d=72$,
$\theta=1$ the ratio per query is $R_I/\Kint^{35/36}$, as in
\cref{thm:comparison}. The finer factor $(\bar\alpha_v/\alpha)^{\theta/d}$ of
the main price is not used: it may not be multiplied into the additive error
of \cref{prop:record:density-structured}.
\end{remark}

\subsection{Energies of the dominating law, and $Y_\lambda\ge(1-L_0-o(1))Y$}

Neither the comparison law nor the bracket measure need be reflection
positive along a sorting sequence, as the structured-energy estimate
requires of the outside law (for the bracket, (CP3) gives positivity at
$\iota$ only).
Their energies are bounded by expanding them into admissible positive laws,
whose energies \cref{prop:record:density-structured} bounds by their prices,
and into record differences grouped in tags; at the cut of a tag all its
differences are applied sharply together
(\cref{lem:record:folded-energy}(d$'$)), and the folded-energy lemma bounds
them by the geometric mean of two folded energies. This is done once
(\cref{lem:sharp-first}); its consequences are
the first comparison, the factor $B$ under the bracket, and the order
relations that Part~V uses throughout.

\paragraph{Energies of selected laws.}
Fix structured slot data: divisors $D_{0,v}$ of $\Mzero$, lengths
$M_v\in\{\Lsc,N\}$ and a law $\mathbf P$ of index vectors
$\omega=(I_v,s_v)_{v\in V}$ independent of the labels, as in
\cref{def:record:structured-energy}(c), with inputs
$g^\omega_v=g^{I_v,D_{0,v},s_v}$. For a finite set $X$ of atoms let $\Pi_X$ be
the coarse price of these data with selected set $X$; for unchanged inputs
($F_A$ at every slot) $\Pi_X=C_X$. An \emph{outside law} is one of the
products $\mu$, $\Nbrz$, $\Nbr$ (with $\Nbrz$ as in
\cref{def:cp3}); write $O_{-X}=\bigotimes_{\Blk\notin X}O_\Blk$.
For a finite signed measure $\kappa$ on $X_X^V$ put
\[
 \mathfrak E_O[\kappa]=\E_{\omega\sim\mathbf P}\int\Lambda_{O_{-X}}
 \bigl((g^\omega_v(z_v,\cdot))_{v\in V}\bigr)\,d\kappa(z),\qquad
 \mathfrak E=\mathfrak E_\mu,
\]
and, for unchanged inputs, $Y^O[\kappa]=Y_{\kappa\otimes O_{-X}}$,
$Y[\kappa]=Y^\mu[\kappa]$. At any cut an energy is the value of the cut form
of the measure at the half-function carrying $F_A$ at every slot of the plus
half, whose reflection is the same half-function \eqref{eq:cc:bilinear}.
Thus $Y[\bigotimes_{\Blk\in X}\mu_\Blk]=Y$, and for admissible positive
$\kappa$ with $q_X\le N^\tau$, $\mathfrak E[\kappa]$ is a structured energy.

\begin{lemma}[Tagged sharp-first expansion]\label{lem:sharp-first}
Assume (CP1) of \cref{thm:comparison}. Let $F$ be a finite set of atoms and
$\mathfrak T$ a finite set of tags with cuts $c_j\in\mathfrak F$, and for
$\Blk\in F$ let $\rho_\Blk$ and $\delta_{\Blk,i}$ ($i$ in a finite set) be
finite signed measures on $X_\Blk^V$ such that:
\begin{enumerate}[label=(\alph*)]
\item $\rho_\Blk-\mu_\Blk=\sum_hc_h\rho_h$ is a finite real combination of
 admissible positive measures, of \emph{variation price}
 $p_\Blk=\sum_h|c_h|\rho_h(1)q_\Blk^{2a}$;
\item $\delta_{\Blk,i}=\kappa'_{\Blk,i}-\kappa_{\Blk,i}$ with admissible positive
 measures of equal mass $w_{\Blk,i}$; it carries a tag $j(\Blk,i)\in\mathfrak T$
 and has a sharp bound $\tilde d_{\Blk,i}$ relative to $\kappa_{\Blk,i}$ at the
 cut $c_{j(\Blk,i)}$ of its tag, in the sense of \cref{def:price}(b): for a
 probability law $\pi$ of shared labels,
 $\kappa_{\Blk,i}=w_{\Blk,i}\int\kappa_s\,\pi(ds)$ and
 $\delta_{\Blk,i}=w_{\Blk,i}\int K_s\,\pi(ds)$, each $K_s$ satisfying
 \eqref{eq:record:sharp} at $c_{j(\Blk,i)}$ relative to the half marginals of
 $\kappa_s$ with the constant $\tilde d_{\Blk,i}$.
\end{enumerate}
Let $d_{\Blk,i}=\tilde d_{\Blk,i}w_{\Blk,i}q_\Blk^{2a}$ be the sharp price of
$\delta_{\Blk,i}$ and $v_{\Blk,i}=2w_{\Blk,i}q_\Blk^{2a}$ the price of its two
positive parts (as in (CP2)(b)), $P=\sum_\Blk p_\Blk$, and let $d_j$ and $v_j$
be the sums of the $d_{\Blk,i}$ and of the $v_{\Blk,i}$ over the differences of
tag $j$. Assume that these prices satisfy the summability condition of
(CP8)(c), as they do in the applications below; in particular $P$ and the
$d_j$, $v_j$ are bounded independently of the scale. Let $X$ be a fixed finite set
of atoms disjoint from $F$ and $\kappa$ an admissible positive measure on
$X_X^V$. Then, for structured slot data and with the selected set $X\cup F$,
\[
 \Bigl|\mathfrak E\Bigl[\kappa\otimes\bigotimes_{\Blk\in F}
 \Bigl(\rho_\Blk+\sum_i\delta_{\Blk,i}\Bigr)\Bigr]
 -\mathfrak E\Bigl[\kappa\otimes\bigotimes_{\Blk\in F}\rho_\Blk\Bigr]\Bigr|
 \le\bigl(e^P\tagser(d,v)+o(1)\bigr)\|\kappa\|\,\Pi_XY ,
\]
with $\tagser$ the series of (CP2). The $o(1)$ consists of energy errors and
one large-product remainder.
\end{lemma}

\begin{proof}
We expand the product over $F$ and bound each nonempty choice of
differences at the cut of its tag; summing over count vectors gives
the series $\tagser(d,v)$. A term of the difference chooses a
nonempty set
$S\subseteq F$ of atoms and one difference $\delta_{\Blk,i_\Blk}$ at each
$\Blk\in S$, and carries $\rho_\Blk$ at the atoms of $F\setminus S$; at most
one difference of an atom enters a term. Its \emph{count vector} is
$\boldsymbol k\in\Z_{\ge0}^{\mathfrak T}\setminus\{0\}$,
$k_j=|\{\Blk\in S:j(\Blk,i_\Blk)=j\}|$. The terms with count vector
$\boldsymbol k$ form the \emph{category} of $\boldsymbol k$. Fix
$\boldsymbol k$ and a tag $j$ with $k_j\ge1$; we bound the sum of the
absolute values of the terms of the category, using the cut $c_j$.

In a term of the category let $S_j\subseteq S$ be the $k_j$ atoms whose
difference has tag $j$. Expand the difference at every atom of
$S\setminus S_j$ into its two given positive parts (not into its Jordan
parts, whose marginals need not be uniform), and every $\rho_\Blk$,
$\Blk\in F\setminus S$, as $\mu_\Blk+\sum_hc_h\rho_h$. The term becomes a
signed sum of energies $T'$ in which the atoms of $S_j$ carry their tag-$j$
differences, a set $X'\subseteq F\setminus S_j$ carries admissible positive
measures, $X$ carries $\kappa$, and every other atom carries $\mu_\Blk$.

In $T'$ apply \cref{lem:record:folded-energy} at $c_j$ with the selected set
$E=X\cup X'\cup S_j$. The outside $\mu_\Blk$ are reflection positive at $c_j$
by (R2) at every cut of the family (CP1). Each $\delta_{\Blk,i_\Blk}$,
$\Blk\in S_j$, is $w_{\Blk,i_\Blk}\int K_s\,\pi_\Blk(ds)$ with kernels that
satisfy \eqref{eq:record:sharp} at the common cut $c_j$; by part (d$'$) of
that lemma, under the product law of the shared labels, their tensor product
over $S_j$ satisfies \eqref{eq:record:sharp} at $c_j$ relative to the
product of the $\kappa_s$, with the constant
$\prod_{\Blk\in S_j}\tilde d_{\Blk,i_\Blk}$. The normalized positive measures on
$X\cup X'$ form a coupling $\pi''$ of its half marginals, so by part (d) the
tensor product with $\pi''$ keeps this constant: positive couplings are
contractions between their marginal spaces, and no positivity of these
factors at $c_j$ is used. By parts (a), (c), $|T'|$ is at most
$\prod_{\Blk\in S_j}\tilde d_{\Blk,i_\Blk}w_{\Blk,i_\Blk}$ times the masses of the
positive factors times the geometric mean of two folded energies
$\mathfrak E_\pm$; by (b) these are energies of admissible positive measures,
since a fold only copies an existing slot, and by (e) their coarse prices
multiply to $\Pi(T')^2$, $\Pi(T')=\Pi_X\prod_{\Blk\in X'\cup S_j}q_\Blk^{2a}$.
If $q_E\le N^\tau$, \cref{prop:record:density-structured}, applied once to
each folded energy with the union $E$ of the selected moduli, gives
\[
 |T'|\le(1+o(1))\prod_{\Blk\in S_j}d_{\Blk,i_\Blk}\prod_{\Blk\in X'}
 \bigl(\text{price of the positive factor at }\Blk\bigr)\,\|\kappa\|\,\Pi_XY :
\]
the geometric mean retains $\Pi(T')$, not $\Pi(T')^2$, no atom is charged
twice, and $1+o(1)$ is not raised to the number of selected atoms. The terms
with $q_E>N^\tau$ have the form $a_E(I^+_EI^-_E)^{1/2}$ of
\cref{lem:record:density-lp}.

Summing over the energies $T'$ of one term: at an atom of $S\setminus S_j$
the two positive parts have total price $v_{\Blk,i_\Blk}$, and at an atom of
$F\setminus S$ the factor is $\mu_\Blk$ (weight one) or a positive record of
$\rho_\Blk-\mu_\Blk$ (total weight $p_\Blk$). Summing then over the terms of
the category, and dropping the requirement that the atoms be distinct, which
only adds nonnegative terms, the category is at most
\[
 (1+o(1))\,\frac{d_j^{k_j}\prod_{j'\ne j}v_{j'}^{k_{j'}}}
 {\prod_{j'\in\mathfrak T}k_{j'}!}\,\prod_{\Blk\in F}(1+p_\Blk)\,
 \|\kappa\|\,\Pi_XY ,
\]
because the sum, over the sets of $k_{j'}$ differences of tag $j'$ at
distinct atoms, of the product of their prices is at most the $k_{j'}$-th
power of the total price of the tag divided by $k_{j'}!$. This holds for
every tag $j$ with $k_j\ge1$, and we use the tag that gives the smallest
bound. The minimum is taken after the whole category has been bounded with
each tag, not inside the sum over its terms, so it needs no relation between
the ratios $d_{\Blk,i}/v_{\Blk,i}$ of the individual differences; the two
bounds refer to the same terms, and their intermediate folded couplings need
not coincide. The categories partition the terms. Summing over
$\boldsymbol k\ne0$ and using $\prod_\Blk(1+p_\Blk)\le e^P$ gives the main
term $e^P\tagser(d,v)\|\kappa\|\Pi_XY$, the tag totals over $F$ being at most
the totals of the statement and $\tagser$ being nondecreasing. The terms with
$q_E>N^\tau$ of all categories form one large-product remainder: such a term
is determined by the role of each atom of $E\setminus X$ (a difference
applied sharply, a positive part of a difference, or a positive record of
$\rho_\Blk-\mu_\Blk$), so for each $E$ their coefficients add up to at most
$\|\kappa\|\,\Pi_X$ times the product over $E\setminus X$ of
$p_\Blk+v_\Blk+d_\Blk$, with $v_\Blk=\sum_iv_{\Blk,i}$, $d_\Blk=\sum_id_{\Blk,i}$.
\end{proof}

\begin{corollary}[Two tags]\label{cor:two-cut-series}
Let $\mathfrak T=\{0,1\}$ and $d_0=d_1\le v_0=v_1$, and write $d$ and $v$ for
the common values (here $d$ is a tag total, not the number of slots) and
$\tagser(d,v)$ for the series of (CP2) at these totals. Then
\[
 \tagser(d,v)=\sum_{k_0+k_1>0}\frac{d^{\max(k_0,k_1)}v^{\min(k_0,k_1)}}
 {k_0!\,k_1!}.
\]
For every integer $n\ge1$ with $d<n+2$ and $dv<(n+2)^2$,
\begin{align*}
 \tagser(d,v)\le\tagser_n^+(d,v):={}&\sum_{i=1}^{n}\Bigl[\frac{2d^i}{i!}
 \sum_{l=0}^{i-1}\frac{v^l}{l!}+\frac{(dv)^i}{(i!)^2}\Bigr]\\
 &+\frac{2e^vd^{n+1}}{(n+1)!\,\bigl(1-d/(n+2)\bigr)}
 +\frac{(dv)^{n+1}}{\bigl((n+1)!\bigr)^2\bigl(1-dv/(n+2)^2\bigr)},
\end{align*}
every term of which is nonnegative and nondecreasing in $d$ and in $v$ where
the two gates hold.
\end{corollary}

\begin{proof}
With $d\le v$, $\min(d^{k_0}v^{k_1},d^{k_1}v^{k_0})=d^{\max(k_0,k_1)}
v^{\min(k_0,k_1)}$, and a count vector with $k_1=0$ (or $k_0=0$) allows one
tag only, with the same value; this is the series. Group its terms by
$i=\max(k_0,k_1)\ge1$: the two terms with $\{k_0,k_1\}=\{i,l\}$, $l<i$, give
$(2d^i/i!)\sum_{l<i}v^l/l!$, and the term $k_0=k_1=i$ gives
$(dv)^i/(i!)^2$. The groups with $i\le n$ form the finite sum. For $i\ge n+1$
bound $\sum_{l<i}v^l/l!$ by $e^v$; the successive ratios of $2d^i/i!$ are
$d/(i+1)\le d/(n+2)$, and those of $(dv)^i/(i!)^2$ are
$dv/(i+1)^2\le dv/(n+2)^2$, so the two geometric series give the two tails.
Every term is a power series in $d$ and $v$ with nonnegative coefficients,
or such a series divided by a positive nonincreasing function of $d$ or of
$dv$, and is therefore nonnegative and nondecreasing in $d$ and in $v$.
\end{proof}

In Part~V, $n=4$ (\cref{def:closing-quantities}). The checked
inequalities $d<6$, $dv<36$ and $d\le v$ are listed in
\cref{app:certificates}; these inequalities justify the displayed
geometric bound for the tail of the two-tag series, which itself converges
for every finite nonnegative $d,v$. At the baseline parameters its value
is the one used in \cref{rem:cp:tags} (no comparison with a
one-sharp bound is needed).

\begin{lemma}[First comparison]\label{lem:record:first-stage}
Assume (CP1), (CP2) and (CP8) of \cref{thm:comparison}. Then
$Y_\lambda\ge(1-L_0-o(1))Y$, with $L_0=e^P\tagser(d,v)$ of (CP2), where the
$o(1)$ consists of energy errors and one large-product remainder.
\end{lemma}

\begin{proof}
\emph{Step 1: $Y_\rho\ge Y$,} where $\rho=\bigotimes_\Blk\rho_\Blk$, $\rho_\Blk=\mu_\Blk$
off the first-stage atoms. At the cut $\vartheta_0$ of (CP2)(a),
$\mu_\Blk\succeq_{\vartheta_0}0$ for every atom by (R2) at every cut of the
family, and $\rho_\Blk\succeq_{\vartheta_0}\mu_\Blk$ on the first-stage atoms.
Tensor products preserve the Loewner order of positive semidefinite matrices
(\cref{lem:record:block-gram}(d)), and the cut matrix of a product measure is
the tensor product of the cut matrices \eqref{eq:cc:kron}; so
$\rho\succeq_{\vartheta_0}\mu$, and evaluation at the half-function carrying
$F_A$ at every slot gives $Y_\rho\ge Y_\mu=Y$ \eqref{eq:cc:bilinear}.

\emph{Step 2: $|Y_\lambda-Y_\rho|\le(L_0+o(1))Y$.} This is
\cref{lem:sharp-first} with $X=\varnothing$, $F$ the first-stage atoms of
$\mathcal A_N$, the $\rho_\Blk$ and the tagged record differences of (CP2),
and unchanged inputs ($\Pi_\varnothing=1$); the tag totals over $\mathcal A_N$
are at most those of (CP2).
\end{proof}

\begin{remark}[Why the first stage is cheap]
Step~1 is an exact relation at one scale; in the construction (CP2)(a) is
the polarization inequality \cref{lem:tr:polarization}
(\cref{subsec:middle-first-stage}).
Step~2 pays only for record differences between laws with the same actual
conditional marginals at the cut of their tag, whose sharp price
(an operator distance times a mass) is smaller than their
variation price by a factor $O(p^{-1/2})$ at the least prime $p$
of the atom; and the differences of one tag are paid together at
their
common cut, so that a term with $k_0$ differences of one tag and $k_1$ of the
other costs the smaller of $d^{k_0}v^{k_1}/(k_0!k_1!)$ and
$v^{k_0}d^{k_1}/(k_0!k_1!)$ instead of one sharp price times the variation of
all the others. At the baseline parameters the two tags have $d=D/2$, about
$2.37\cdot10^{-5}$, and $v=V/2$, about $1.69$, and
$L_0<4.76332373\cdot10^{-4}$; the one-sharp cost $De^{P+V}$ would be about
$7.6\cdot10^{-3}$, larger than the margin of the closing by eight orders of
magnitude (\cref{tab:closing-values}).
\end{remark}

\begin{lemma}[Energies under the bracket]\label{lem:bracket-energy}
Assume (CP1), (CP2)(b), (CP3) and (CP8) of \cref{thm:comparison}. Let $X$ be a
fixed finite set of atoms, $\kappa$ an admissible positive measure on $X_X^V$, and
fix structured slot data. Then
\begin{gather*}
 |\mathfrak E[\kappa]|\le(1+o(1))\|\kappa\|\,\Pi_XY,\qquad
 |\mathfrak E_{\Nbrz}[\kappa]|\le(e^{U_0}+o(1))\|\kappa\|\,\Pi_XY,\\
 |\mathfrak E_{\Nbr}[\kappa]-\mathfrak E_{\Nbrz}[\kappa]|
 \le(R_{\rm err}+o(1))\|\kappa\|\,\Pi_XY,
\end{gather*}
and hence $|\mathfrak E_{\Nbr}[\kappa]|\le(B+o(1))\|\kappa\|\,\Pi_XY$. When such
an energy is one of the two factors of a Cauchy--Schwarz step, the same bound
holds for both, and their geometric mean retains $B$, not $B^2$. The $o(1)$
consists of energy errors and large-product remainders.
\end{lemma}

\begin{proof}
For fixed $X$, eventually $q_X\le N^\tau$, and the first bound is
\cref{prop:record:density-structured}. In a sum over $X$, terms with
$q_X>N^\tau$ are retained in the aggregate large-product remainder
described above. On a first-stage atom $\Blk\notin X$
write $\Nbr_\Blk=\Nbr_{0,\Blk}+\delta_{E,\Blk}$ and
$\Nbr_{0,\Blk}=\mu_\Blk+(\Nbr_{0,\Blk}-\mu_\Blk)$, the second summand being an
admissible positive measure of price $u_\Blk$ by (CP3) and (CP8). Expanding
the latter, $\mathfrak E_{\Nbrz}[\kappa]$ is the sum over the sets $X'$ of
first-stage atoms outside $X$ of the energies with selected measures
$\kappa\otimes\bigotimes_{\Blk\in X'}(\Nbr_{0,\Blk}-\mu_\Blk)$, each at most
$(1+o(1))\|\kappa\|\,\Pi_X\prod_{X'}u_\Blk\,Y$; the sum is at most
$(1+o(1))\prod_\Blk(1+u_\Blk)\|\kappa\|\,\Pi_XY\le(1+o(1))e^{U_0}\|\kappa\|\,\Pi_XY$.
The difference is \cref{lem:sharp-first} with $F$ the first-stage atoms
outside $X$, $\rho_\Blk=\Nbr_{0,\Blk}$ (variation price $u_\Blk$) and the tagged
record differences of (CP2)(b): it is at most
$(e^{U_0}\tagser(d,v)+o(1))\|\kappa\|\,\Pi_XY=(R_{\rm err}+o(1))\|\kappa\|\,
\Pi_XY$, the tag totals over the first-stage atoms outside $X$ being at most
those of (CP2) and $\tagser$ being nondecreasing. For
the last assertion, the two factors of a Cauchy--Schwarz step have coarse
prices whose product is the square of the original price
(\cref{lem:record:folded-energy}(e)).
\end{proof}

The next lemma collects the order, contraction and interpolation
estimates for energies that Part~V applies repeatedly: its parts
(E1)--(E3) are the basic comparisons, and (P1)--(P5), (P3$'$) are
the product estimates used when a bound proved in a fixed chart is
transported to actual marginals (restoration), when centred norms
are interpolated, and when class bounds are formed, in the places
named in Part~V.

\begin{lemma}[Energy toolkit: order, contraction and interpolation estimates]\label{lem:energy-toolkit}
Assume (CP1)--(CP3) and (CP8) of \cref{thm:comparison}. Let $X$ be a finite
set of atoms and $\kappa,\kappa'$ finite signed measures on $X_X^V$. In (E2),
(P1)--(P3) and (P3$'$), $\vartheta$ is a cut of the family and $O=\mu$, or
$\vartheta=\iota$ and $O$ is any outside law.
\begin{enumerate}[label=(E\arabic*)]
\item\label{it:record:E1} If $\kappa$ is positive and admissible, then
 $|Y[\kappa]|\le(1+o(1))\|\kappa\|C_XY$,
 $|Y^{\Nbrz}[\kappa]|\le(e^{U_0}+o(1))\|\kappa\|C_XY$,
 $|Y^{\Nbr}[\kappa]|\le(B+o(1))\|\kappa\|C_XY$ and
 $|Y^{\Nbr}[\kappa]-Y^{\Nbrz}[\kappa]|\le(R_{\rm err}+o(1))\|\kappa\|C_XY$.
\item\label{it:record:E2} If $\kappa\succeq_\vartheta\kappa'$, then
 $Y^O[\kappa]\ge Y^O[\kappa']$; more generally, at any cut $\vartheta$,
 $\kappa\succeq_\vartheta\kappa'$ implies $Y_{\kappa\otimes O'}\ge Y_{\kappa'\otimes O'}$
 for every finite signed measure $O'$ on the labels of the atoms outside $X$
 with $O'\succeq_\vartheta0$. If $\kappa\succeq_\iota0$, then
 $Y[\kappa]\le Y^{\Nbrz}[\kappa]$.
\item\label{it:record:E3} A positive measure fixed samplewise by the
 reflection of $\vartheta$ (\emph{diagonal at $\vartheta$}) has a
 nonnegative diagonal cut matrix at $\vartheta$, so it is $\succeq_\vartheta0$.
\end{enumerate}
For a positive measure $m$ on the plus-half labels at $\vartheta$ let
$\operatorname{Diag}_\vartheta(m)$ be its diagonal copy of
\cref{lem:commoncut}(i) (the plus half has law $m$, the minus half is its
reflection) and $J_\vartheta(m)=(m\otimes m)/m(1)$ the measure with
independent halves of laws proportional to $m$.
\begin{enumerate}[label=(P\arabic*)]
\item\label{it:record:P1} \emph{Contraction.} If $\pi$ is a positive measure on
 $X_X^V$ with half marginals $m_\pm$ at $\vartheta$, then
 $|Y^O[\pi]|\le Y^O[\operatorname{Diag}_\vartheta(m_+)]^{1/2}
 Y^O[\operatorname{Diag}_\vartheta(m_-)]^{1/2}$, where
 $\operatorname{Diag}_\vartheta(m_-)$ is the diagonal measure whose minus half
 has law $m_-$.
\item\label{it:record:P2} \emph{Density domination.} If $m\le Mm'$ for positive
 measures on the plus-half labels, then
 $Y^O[\operatorname{Diag}_\vartheta(m)]\le MY^O[\operatorname{Diag}_\vartheta(m')]$.
\item\label{it:record:P3} \emph{Centred interpolation.} Let $\pi$ be a positive
 measure symmetric at $\vartheta$ with half marginal $m$ and Markov operator
 $T$ on $L^2(m)$. If $\|T-\Jproj\|\le\kappa_0\le1$, then
 $\pi\preceq_\vartheta(1-\kappa_0)J_\vartheta(m)
 +\kappa_0\operatorname{Diag}_\vartheta(m)$, hence $Y^O[\pi]\le(1-\kappa_0)
 Y^O[J_\vartheta(m)]+\kappa_0Y^O[\operatorname{Diag}_\vartheta(m)]$. The case
 $\kappa_0=1$ holds for every such $\pi$.
\item[(P3$'$)]\setprimedlabel{(P3$'$)}\label{it:record:P3b}
 \emph{Product centred interpolation.} Let $\pi$ be as in (P3), on the labels
 of $X$, with half marginal $m$ and constant $\kappa_0\le1$, and let $\pi'$ be
 a positive measure on the labels of a finite set $Z$ of atoms disjoint from
 $X$, symmetric at $\vartheta$, with half marginal $m'$. Then
 \[
  \pi\otimes\pi'\preceq_\vartheta(1-\kappa_0)\,J_\vartheta(m)\otimes
  \operatorname{Diag}_\vartheta(m')+\kappa_0\operatorname{Diag}_\vartheta(m\otimes m'),
 \]
 hence $Y^O[\pi\otimes\pi']\le(1-\kappa_0)Y^O[J_\vartheta(m)\otimes
 \operatorname{Diag}_\vartheta(m')]+\kappa_0Y^O[\operatorname{Diag}_\vartheta(m\otimes
 m')]$. No order is tensored with an indefinite factor: $\pi'$ enters only
 through its Markov operator, a self-adjoint contraction.
\item\label{it:record:P4} \emph{Products.} If $\pi_\Blk$ $(\Blk\in X)$ are as in
 (P3) with half marginals $m_\Blk$ and constants $\kappa_\Blk$, then
 $\bigotimes_\Blk\pi_\Blk$ is as in (P3) with half marginal
 $\bigotimes_\Blk m_\Blk$ and constant $\max_\Blk\kappa_\Blk$.
\item\label{it:record:P5} \emph{$\pm$ tensorization.} If $\zeta$ is a signed
 measure on $X_X^V$ with $\kappa\pm\zeta\succeq_\iota0$, then
 $|Y_{\zeta\otimes\lambda_{-X}}|\le Y^{\Nbr}[\kappa]$. If
 $\bigl(\begin{smallmatrix}\kappa&\zeta\\ \zeta^*&\kappa'\end{smallmatrix}
 \bigr)\succeq_\iota0$, then $|Y_{\zeta\otimes\lambda_{-X}}|\le
 Y^{\Nbr}[\kappa]^{1/2}Y^{\Nbr}[\kappa']^{1/2}$. If the cut matrix of $\zeta$ at
 $\iota$ is skew, then $Y_{\zeta\otimes\lambda_{-X}}=0$.
\end{enumerate}
\end{lemma}

\begin{proof}
\ref{it:record:E1} is \cref{lem:bracket-energy} for unchanged inputs.
\ref{it:record:E2}: for $O'$ as stated, the cut matrix of
$(\kappa-\kappa')\otimes O'$ at $\vartheta$ is
$C_\vartheta(\kappa-\kappa')\otimes C_\vartheta(O')$ \eqref{eq:cc:kron}, a
Kronecker product of symmetric positive semidefinite matrices; evaluate at the
half-function carrying $F_A$ \eqref{eq:cc:bilinear}. The outside laws are such
measures: $\mu_{-X}\succeq_\vartheta0$ at every cut of the family by (R2), as
tensor products of reflection positive measures are reflection positive
(\cref{def:rp}); $\Nbr_{-X}\succeq_\iota0$ because each $\Nbr_\Blk$ is a
diagonal block of the bracket of (CP3); and $\Nbr_{0,-X}\succeq_\iota0$
because $\Nbr_{0,\Blk}=\mu_\Blk+(\Nbr_{0,\Blk}-\mu_\Blk)$, with
$\mu_\Blk\succeq_\iota0$ and the second summand positive and reserved
diagonal, hence $\succeq_\iota0$ by (E3). For the last claim, $0\preceq_\iota\mu_\Blk\preceq_\iota
\Nbr_{0,\Blk}$, so $\mu_{-X}\preceq_\iota\Nbr_{0,-X}$ (\cref{lem:record:block-gram}(d)),
and $C_\iota(\kappa)\otimes(C_\iota(\Nbr_{0,-X})-C_\iota(\mu_{-X}))\succeq0$.
\ref{it:record:E3} is \cref{lem:commoncut}(i).

\ref{it:record:P1} Let $D_\pm$ be the diagonal cut matrices of
$\operatorname{Diag}_\vartheta(m_\pm)$, the minus half read on the plus half
through the reflection. For real $f,g$ on the plus-half labels,
$\ip f{D_+f}+2\ip f{C_\vartheta(\pi)g}+\ip g{D_-g}=\int(f(z_+)+g(\tilde z_-))^2
\,d\pi\ge0$, with $\tilde z_-$ the reflected minus half, so
$\bigl(\begin{smallmatrix}D_+&C_\vartheta(\pi)\\
C_\vartheta(\pi)^{\mathsf T}&D_-\end{smallmatrix}\bigr)\succeq0$. Tensorize with
$\bigl(\begin{smallmatrix}1&1\\1&1\end{smallmatrix}\bigr)\otimes
C_\vartheta(O_{-X})$ and apply the block-Gram inequality
(\cref{lem:record:block-gram}) at the half-function carrying $F_A$; for
$O=\mu$ this is \cref{lem:record:folded-energy}(a) with $d=1$.
\ref{it:record:P2}: $\operatorname{Diag}_\vartheta(Mm'-m)$ is positive and
diagonal at $\vartheta$; apply (E3), (E2).

\ref{it:record:P3} Relative to $m$ the cut matrices of $\pi$,
$J_\vartheta(m)$ and $\operatorname{Diag}_\vartheta(m)$ are $D_mT$, $D_m\Jproj$
and $D_m$ ($D_m$ the diagonal matrix of $m$), so the cut order between
measures of half marginal $m$ is the Loewner order of their operators on
$L^2(m)$. $T$ is self-adjoint and fixes the constants, so
$T-\Jproj=(I-\Jproj)(T-\Jproj)(I-\Jproj)\preceq\kappa_0(I-\Jproj)$, that is,
$T\preceq(1-\kappa_0)\Jproj+\kappa_0I$; apply (E2). The Markov operator of a
positive coupling is a contraction on $L^2(m)$ (Jensen), so
$\|T-\Jproj\|\le1$ always.

\ref{it:record:P3b} Relative to $m\otimes m'$ the cut matrices of
$\pi\otimes\pi'$, $J_\vartheta(m)\otimes\operatorname{Diag}_\vartheta(m')$ and
$\operatorname{Diag}_\vartheta(m\otimes m')$ are $D(T\otimes T')$,
$D(\Jproj\otimes I)$ and $D$, with $D$ the diagonal matrix of $m\otimes m'$
and $T'$ the Markov operator of $\pi'$ on $L^2(m')$, which is self-adjoint,
because $\pi'$ is symmetric, and a contraction. The operator $T\otimes T'$
maps $\operatorname{ran}\Jproj\otimes L^2(m')$ and its orthogonal complement
$L^2_0(m)\otimes L^2(m')$ into themselves. On the first it is $I\otimes T'\preceq
I$; on the second it is $(T-\Jproj)\otimes T'$, self-adjoint of norm at most
$\kappa_0$, hence $\preceq\kappa_0I$. So
$T\otimes T'\preceq\Jproj\otimes I+\kappa_0(I-\Jproj)\otimes I
=(1-\kappa_0)\Jproj\otimes I+\kappa_0I$, which is the order claimed; apply
(E2).

\ref{it:record:P4}: the Markov operator of the
product is $\bigotimes_\Blk T_\Blk$, self-adjoint, and
$\|\bigotimes_\Blk T_\Blk-\Jproj\|=\max_\Blk\|T_\Blk-\Jproj\|$ (\cref{lem:sector}(c)).

\ref{it:record:P5} By (CP3), $\Nbr_\Blk\pm\lambda_\Blk\succeq_\iota0$ for every
atom (off the first-stage atoms these are $2\mu_\Blk$ and $0$). With
$\kappa\pm\zeta\succeq_\iota0$, \cref{lem:record:block-gram}(d) gives
$\kappa\otimes\Nbr_{-X}\pm\zeta\otimes\lambda_{-X}\succeq_\iota0$, and evaluation at
the half-function carrying $F_A$ gives the first claim. For the bracket form
tensorize the bracket of $\zeta$ with those of (CP3)
(\cref{lem:record:block-gram}(b)) and apply part (a) of that lemma. If
$C_\iota(\zeta)$ is skew, so is $C_\iota(\zeta)\otimes C_\iota(\lambda_{-X})$, whose
form vanishes at two equal arguments (\cref{rem:cut-symmetries}(a)).
\end{proof}

\subsection{Expanding $\nu-\lambda$ over changed blocks: the bound as a power series in prices}

The difference $Y_\nu-Y_\lambda$ at the reserved cut is expanded
by choosing, at any set of atoms, which records were changed; each
such choice is a \emph{selection}. A selection in no special class
is bounded by $B$ times the product of the prices of its changed
records. Selections in a class of (CP5) may have a smaller proved
bound. Summing over all selections gives a power series in the
prices; any series with nonnegative coefficients that dominates
it coefficient by coefficient (a \emph{complete polynomial},
\cref{subsec:intro:glossary}) may be used in the closing. This
subsection makes the expansion precise.

\paragraph{Selections.}
Recall from (CP4)(b), (c) (\cref{def:record:comparison-law}) the decomposition
$\nu_\Blk-\lambda_\Blk=\sum_{t\in\mathcal T_\Blk}(\zeta^{\rm s}_{\Blk,t}+
\zeta^{\rm a}_{\Blk,t})$ over the changed records $t$ of $\Blk$, the brackets of
$\zeta^\epsilon_{\Blk,t}$ with the norm envelopes
$\tilde\pi^\epsilon_{\Blk,t}D^\epsilon_{\Blk,t}$, the prices
$\pi^\epsilon_{\Blk,t}$ with totals $S$ and $A$, and the selections $\omega$ (a
set $X$ of atoms with a record $t_\Blk$ and a type $\epsilon_\Blk$ at each, an
even number of them skew) with their measures $\zeta_\omega$, monomials
$\pi_\omega$ and terms $I_\omega=Y_{\zeta_\omega\otimes\lambda_{-X}}$.
In formal variables $x=(x^\epsilon_{\Blk,t})$ for the prices and
$y=(y_{\Blk,j})$ for the auxiliary prices of (CP5) let
$S(x)=\sum x^{\rm s}_{\Blk,t}$, $A(x)=\sum x^{\rm a}_{\Blk,t}$ and
$e^{S(x)}\cosh A(x)-1=\sum_{\boldsymbol n}c_{\boldsymbol n}x^{\boldsymbol n}$, the
sum over the nonconstant monomials of even skew degree. Every
$c_{\boldsymbol n}$ is positive (it is $1/\boldsymbol n!$), and a squarefree
monomial has coefficient one; so has, in particular, the monomial
$x^{\boldsymbol n(\omega)}$ of a selection, for which
$\pi^{\boldsymbol n(\omega)}=\pi_\omega$. A selection determines its atoms,
records and types, hence its monomial; selections of different patterns of
types, also those without skew types, are different monomials, and (CP5) may
be proved for every pattern with an even number of skew types. For a
selection of a class $\mathcal C$ of (CP5) with a class cost let
$c_\omega=c_{\mathcal C}$, and $c_\omega=B$ for a selection in no class; for
these put $b_\omega=c_\omega x^{\boldsymbol n(\omega)}$. For a selection of a
class with class bounds put $b_\omega=\phi_\omega$, a polynomial in $y$. A
monomial in $x$ is the monomial of at most one selection, and a monomial of a
class bound determines its selection (CP5); so $\sum_\omega b_\omega$, the sum
over all selections of atoms of $\mathcal A$, is a power series in $(x,y)$
with nonnegative coefficients, each of whose coefficients is a finite sum
of coefficients of the $b_\omega$.

\begin{lemma}[Record expansion and complete polynomial]
\label{lem:record:record-expansion}
Assume (CP0)--(CP4) and (CP8) of \cref{thm:comparison}.
\begin{enumerate}[label=(\roman*)]
\item $Y_\nu-Y_\lambda=\sum_\omega I_\omega$, and $|I_\omega|\le(B+o(1))\pi_\omega Y$
 for every fixed selection $\omega$.
\item \emph{Complete polynomials.} Assume (CP5). Let $\Psi(x,y)$ be a power
 series with nonnegative coefficients such that $\Psi-\sum_\omega b_\omega$ has
 nonnegative coefficients. Then $|Y_\nu-Y_\lambda|\le(\Psi(\hat\pi,\hat
 y)+o(1))Y$ for all upper prices $\hat\pi\ge\pi$ and upper auxiliary prices
 $\hat y\ge y$ at which $\Psi$ converges. Such a $\Psi$ is a \emph{complete
 polynomial}. Without classes, $\Psi=B(e^{S(x)}\cosh A(x)-1)$, and this
 reads $|Y_\nu-Y_\lambda|\le(B(e^S\cosh A-1)+o(1))Y$; with $B'\ge B$ (for
 instance a certified upper bound for $B$) every $B$ may be replaced by
 $B'$. For classes with class costs, $\Psi=\sum_{\boldsymbol n}(B'-\hat
 g_{\boldsymbol n})c_{\boldsymbol n}x^{\boldsymbol n}$ is complete whenever
 $0\le\hat g_{\boldsymbol n}\le B'$ for every $\boldsymbol n$ and $\hat
 g_{\boldsymbol n(\omega)}\le B'-c_\omega$ for every selection $\omega$ (the
 $\hat g_{\boldsymbol n}$ are arbitrary in $[0,B']$ on the other monomials,
 which have a repeated atom).
\item $\Psi$ is nondecreasing in each price variable and in each auxiliary
 variable. In particular an upper bound for the total of a nonempty set of variables
 may be distributed over them in proportion to their actual values when
 their total is positive, and arbitrarily as nonnegative upper prices when
 their total is zero, giving a sum exactly equal to the given bound.
\item \emph{Completion.} Assume (CP5) with class costs, and let $B'\ge B$. Let
 $\mathcal D(x)=\sum_{\boldsymbol n}d_{\boldsymbol n}x^{\boldsymbol n}$ be a power
 series with nonnegative coefficients, convergent at $\hat\pi$, with
 $d_{\boldsymbol n}\le B'-c_\omega$ on the monomial of every selection and
 $d_{\boldsymbol n}\le B'c_{\boldsymbol n}$ on every other monomial
 ($c_{\boldsymbol n}=0$ if $x^{\boldsymbol n}$ does not occur). Then, with
 $\hat S=S(\hat\pi)$, $\hat A=A(\hat\pi)$,
 \[
  |Y_\nu-Y_\lambda|\le\bigl(B'(e^{\hat S}\cosh\hat A-1)-\mathcal D(\hat\pi)
  +o(1)\bigr)Y .
 \]
 That is, $\Psi=B'(e^{S(x)}\cosh A(x)-1)-\mathcal D(x)$ is a complete
 polynomial in the sense of (ii).
\item \emph{Degree-class replacement.} Divide the skew records into kinds
 $1,\dots,r$: for every kind $l$ a set $\mathcal R_l$ of pairs $(\Blk,t)$,
 $t\in\mathcal T_\Blk$, the sets being disjoint. Let
 $A_l(x)=\sum_{(\Blk,t)\in\mathcal R_l}x^{\rm a}_{\Blk,t}$, and let every
 $(\Blk,t)\in\mathcal R_l$ carry auxiliary prices
 $y^{\rm arm}_{\Blk,t},y^{\rm lead}_{\Blk,t},y^{\rm rem}_{\Blk,t}\ge0$ (raw-arm,
 leading and remainder), attached to $(t,{\rm a})$, with totals
 $A^{\rm arm}_l$, $A^{\rm lead}_l$, $A^{\rm rem}_l$ over $\mathcal R_l$. For
 $\boldsymbol n\in\Z_{\ge0}^r$ with $|\boldsymbol n|\ge2$ even let $\mathcal
 C_{\boldsymbol n}$ be the class of the selections consisting of $n_l$ skew
 records of kind $l$ for every $l$ and of no other record, and suppose that
 $\mathcal C_{\boldsymbol n}$ satisfies (CP5) with
 \[
  \phi_\omega=C_{\boldsymbol n}\prod_{\Blk\in X}y^{\rm arm}_{\Blk,t_\Blk}
  +B'\Bigl(\prod_{\Blk\in X}\bigl(y^{\rm lead}_{\Blk,t_\Blk}
  +y^{\rm rem}_{\Blk,t_\Blk}\bigr)-\prod_{\Blk\in X}y^{\rm lead}_{\Blk,t_\Blk}\Bigr)
 \]
 for a number $C_{\boldsymbol n}\ge0$, where $X=X(\omega)$. Let $\Psi_0$ be a
 complete polynomial for the other classes, the selections of $\mathcal
 C_{\boldsymbol n}$ being in no class, whose coefficients on the monomials of
 degree $n_l$ in the skew variables of kind $l$ for every $l$ and of degree
 zero in all other variables are those of $B'(e^{S(x)}\cosh A(x)-1)$. Then,
 with $\boldsymbol n!=\prod_ln_l!$,
 \[
  \Psi=\Psi_0-\frac1{\boldsymbol n!}\Bigl[B'\prod_lA_l^{n_l}
  -C_{\boldsymbol n}\prod_l\bigl(A^{\rm arm}_l\bigr)^{n_l}
  -B'\Bigl(\prod_l\bigl(A^{\rm lead}_l+A^{\rm rem}_l\bigr)^{n_l}
  -\prod_l\bigl(A^{\rm lead}_l\bigr)^{n_l}\Bigr)\Bigr],
 \]
 read as a power series in the variables of the individual records, is a
 complete polynomial for all the classes. For a set of such degree vectors
 the replacements add. \emph{Credits:} if moreover the selections $\omega$ of
 $\mathcal C_{\boldsymbol n}$ satisfy (CP5) with the smaller class bounds
 $\phi_\omega-\gamma_\omega$, where the credits $\gamma_\omega$ are polynomials
 with nonnegative coefficients (for instance variance credits on the raw-arm
 monomials), then $\Psi-\Gamma$ is a complete polynomial for every power
 series $\Gamma$ with nonnegative coefficients whose coefficient is at most
 that of $\gamma_\omega$ on every monomial of a selection $\omega$ of $\mathcal
 C_{\boldsymbol n}$ and at most that of $\Psi-\sum_\omega b_\omega$ on every
 other monomial, as in (iv).
\end{enumerate}
The $o(1)$ consists of energy errors and large-product remainders.
\end{lemma}

In (ii) the bounds of (CP5) hold robustly and uniformly in the class
((CP8)(b)); this is where the bounded arity $n_{\max}$ enters, since the
bounds of the classes then come from finitely many types of operator
reductions. No argument uses a particular arity. In Part~V every selection of
the variance, background and late-pool credits and of the envelope classes
has at most six atoms, and so has every selection of the mean savings of
(CP7) (\cref{prop:background-variance,lem:pool-coupons}).

The replacement (v) applies to coarser kinds as well: if $C_{\boldsymbol n}$
depends only on the sums of the $n_l$ over groups of kinds, then by the
multinomial expansion the replacements of the degree vectors with given group
sums add up to the replacement for the merged kinds, whose totals are the
sums of the totals. In Part~V the kinds are the fold records and the two
matching records $E_0$, $E_1$: the general classes, with $n$ fold and $m$
matching records, are replaced with the merged kind of the matching records,
and the matching pairs with $E_0$ and $E_1$ separate, the aligned pairs (two
records of one kind) and the crossed pairs (one of each) each having the
total $A_E^2/4$, the two matching records having equal totals
(\cref{cor:envelope-costs,def:closing-quantities}).

\begin{proof}
(i) Expand $\nu=\bigotimes_\Blk(\lambda_\Blk+\sum_t(\zeta^{\rm s}_{\Blk,t}+\zeta^{\rm a}_{\Blk,t}))$.
A term with an odd number of skew factors has a skew cut matrix at $\iota$,
because $\lambda_\Blk$ is symmetric there by (CP3), and vanishes by
\ref{it:record:P5}. The remaining terms are $Y_\lambda$ and the $I_\omega$.
For a selection, the tensor product of the brackets of its factors and
\ref{it:record:P5} give $|I_\omega|\le\prod_{\Blk\in X}
\tilde\pi^{\epsilon_\Blk}_{\Blk,t_\Blk}\,Y^{\Nbr}[\bigotimes_{\Blk\in X}
D^{\epsilon_\Blk}_{\Blk,t_\Blk}]\le(B+o(1))\pi_\omega Y$ by \ref{it:record:E1}, the
laws $D^\epsilon_{\Blk,t}$ being admissible and normalized.

(ii) Fully expand the bracket outside each selection as in the proof of
\cref{lem:bracket-energy}. For selections in no class, apply the uniform
structured-energy bound only to the resulting terms whose full union
$E$ satisfies $q_E\le N^\tau$; the other terms, summed over the selections
and the outside expansion together, are one large-product remainder.
Their coefficients are bounded by products of the sums of the local
prices in (CP8)(c), so \cref{lem:record:density-lp} makes their total
$o(1)Y$. The small terms have one uniform relative energy error.
Combine these bounds with the class sums of (CP5), which are uniform
by (CP8)(b). Thus, after adding this aggregate remainder,
$|Y_\nu-Y_\lambda|\le\sum_\omega|I_\omega|\le\sum_\omega\bigl(b_\omega(\pi,y)
+o(1)\pi_\omega\bigr)Y+o(1)Y$, with $b_\omega(\pi,y)=c_\omega\pi_\omega$ off the classes
with class bounds and $b_\omega(\pi,y)=\phi_\omega(y)$ on them, the bounds of
(CP5) being uniform in the class ((CP8)(b)). The arguments $(\pi,y)$ and
$(\hat\pi,\hat y)$ are nonnegative and $(\pi,y)\le(\hat\pi,\hat y)$, so
$\sum_\omega b_\omega(\pi,y)\le\Psi(\pi,y)\le\Psi(\hat\pi,\hat y)$ by the
dominance and by (iii). Also $\sum_\omega\pi_\omega\le e^S\cosh A-1$ is
bounded, so the errors add up to $o(1)Y$. Without classes
$\sum_\omega b_\omega=B\sum_\omega x^{\boldsymbol n(\omega)}$, which
$B(e^{S(x)}\cosh A(x)-1)$ dominates, its coefficient being $B$ on every
squarefree monomial and nonnegative on the others. For the last
sentence of (ii): the coefficient of $\Psi$ on the monomial of a selection is
$B'-\hat g_{\boldsymbol n(\omega)}\ge c_\omega$, since
$c_{\boldsymbol n(\omega)}=1$, and on the other monomials it is nonnegative;
distinct selections have distinct monomials, so this is the dominance.

(iii) Every coefficient of $\Psi$ is nonnegative. If
$0<\sum_G\pi\le\hat\Pi_G$, then
$\hat\pi_g=\pi_g\hat\Pi_G/\sum_G\pi\ge\pi_g$ and
$\sum_G\hat\pi_g=\hat\Pi_G$. If $G$ is nonempty and the actual total
is zero, all its prices are zero; assign the upper total to any one variable
and zero to the others. Empty groups need no allocation.

(iv) On the monomial of a selection the coefficient of
$B'(e^{S(x)}\cosh A(x)-1)-\mathcal D(x)$ is $B'-d_{\boldsymbol n}\ge c_\omega$,
and on every other monomial it is $B'c_{\boldsymbol n}-d_{\boldsymbol n}\ge0$;
this is the dominance of (ii), both series converging absolutely at $\hat\pi$.

(v) By the multinomial expansion of $\cosh A(x)$, the monomials of
$B'(e^{S(x)}\cosh A(x)-1)$ of degree $n_l$ in the skew variables of kind $l$
for every $l$ and of degree zero in all other variables sum to
$B'\prod_lA_l(x)^{n_l}/\boldsymbol n!$: the coefficient of such a monomial
$x^{\boldsymbol m}$ in $\prod_lA_l^{n_l}$ is $\boldsymbol n!/\boldsymbol m!$,
and in $e^{S}\cosh A$ it is $1/\boldsymbol m!$. These monomials are the
monomials of the selections of $\mathcal C_{\boldsymbol n}$ and of the tuples
of that degree with a repeated atom. By the hypothesis on $\Psi_0$,
subtracting $B'\prod_lA_l^{n_l}/\boldsymbol n!$ sets their coefficients to
zero and leaves all other coefficients of $\Psi_0$, which dominate the
$b_\omega$ of the selections outside $\mathcal C_{\boldsymbol n}$ (those of
$\mathcal C_{\boldsymbol n}$, in no class for $\Psi_0$, contributed only to
these monomials). The added polynomial
$\bigl[C_{\boldsymbol n}\prod_l(A^{\rm arm}_l)^{n_l}+B'\bigl(\prod_l(A^{\rm
lead}_l+A^{\rm rem}_l)^{n_l}-\prod_l(A^{\rm lead}_l)^{n_l}\bigr)\bigr]/\boldsymbol
n!$ is a polynomial in the auxiliary variables with nonnegative coefficients,
so $\Psi$ has nonnegative coefficients. It dominates
$\sum_{\omega\in\mathcal C_{\boldsymbol n}}\phi_\omega$. A selection of
$\mathcal C_{\boldsymbol n}$ is a set of distinct atoms with one record of
kind $l$ at $n_l$ of them; it corresponds to exactly $\boldsymbol n!$ tuples
$((\Blk_1,t_1),\dots,(\Blk_{|\boldsymbol n|},t_{|\boldsymbol n|}))$ whose first
$n_1$ entries lie in $\mathcal R_1$, the next $n_2$ in $\mathcal R_2$, and so
on, with distinct atoms. So $\sum_{\omega\in\mathcal C_{\boldsymbol
n}}\phi_\omega$ is $1/\boldsymbol n!$ times the sum of
$C_{\boldsymbol n}\prod_iy^{\rm arm}_{\Blk_i,t_i}+B'\bigl(\prod_i(y^{\rm
lead}_{\Blk_i,t_i}+y^{\rm rem}_{\Blk_i,t_i})-\prod_iy^{\rm lead}_{\Blk_i,t_i}\bigr)$
over these tuples. The sum of the same expression over all tuples, the atoms
not required to be distinct, is the added polynomial, by expanding the
products of the totals; and every expression in the sum has nonnegative
coefficients, the bracket being the sum of the products in which at least
one remainder price occurs. So dropping the distinctness adds only
nonnegative coefficients, and $\Psi$ dominates $\sum_\omega b_\omega$. Several
degree vectors concern disjoint sets of monomials in $x$ and add disjoint
nonnegative polynomials in $y$, so their replacements add. For the credits:
on a monomial of a selection $\omega$ of
$\mathcal C_{\boldsymbol n}$ the coefficient of $\Psi-\Gamma$ is at least
that of $\phi_\omega-\gamma_\omega$, and on every other monomial at least that
of $\sum_\omega b_\omega$; this is the dominance of (ii) for the credited
class bounds, and it implies that $\Psi-\Gamma$ has nonnegative coefficients.
\end{proof}

\begin{remark}[Completion and subtraction]\label{rem:completion-subtraction}
Parts (iv) and (v) are the only ways in which a discount is subtracted from
the reserved cost, and they fix the order of operations of the closing
(\cref{cert:closing}, which takes for $B'$ the certified upper end of $B$):
every genuine selection first receives its nonnegative coefficient, the cost
of its class after its credits ($B'-d_{\boldsymbol n}\ge c_\omega$ in (iv)) or
its credited class bound (in (v)); the monomials with a repeated atom receive
nonnegative coefficients ($[0,B'c_{\boldsymbol n}]$ in (iv)); only then is the
series rewritten as an exponential minus homogeneous polynomials in the
totals of the price variables and of the auxiliary variables and evaluated at
the upper prices. The homogeneous terms are not lower bounds for a
distinct-block discount, and no claim is made that a subtracted polynomial,
considered alone, is a lower bound for the saving from distinct blocks.
The finite gates of the closing which say that every coefficient
stays nonnegative after the credits (every leading class cost $C_{n,m}\ge0$,
every credited coefficient nonnegative, every restoration gain
--- the factor in $[0,1)$ by which a restored variance credit
(\cref{sec:variance}) multiplies a class coefficient --- in
$[0,1)$) are
exactly the conditions of (iv) and (v); that a class is replaced only when
its saving is positive is a choice, not a condition of (v), which holds
whatever the sign of the saving. Two rules govern the subtractions in this
paper: an upper bound for a
quantity is never used as a lower bound for it; and a discount is subtracted
only inside an expression with nonnegative coefficients in the same chosen
variables, as in (iv) and (v). In the replaced form (v) the remainder
polynomial contains only the separately proved leading prices $y^{\rm lead}$
and the remainder prices $y^{\rm rem}$; neither a raw arm price nor the
difference of a full price and a remainder price may replace a leading price
there, since neither bounds the leading part of a record change.

The same completion bounds other polynomials with nonnegative coefficients.
If $0\le\kappa_b\le\kappa_H$ and $v_\Blk,w_\Blk\ge0$, then
\[
 \kappa_H\prod_\Blk(1+v_\Blk+w_\Blk)-(\kappa_H-\kappa_b)\prod_\Blk(1+v_\Blk)
\]
has coefficient $\kappa_b$ on the monomials in the $v_\Blk$ alone and $\kappa_H$
on all others, all squarefree, and it is nondecreasing in every $v_\Blk$ and
every $w_\Blk$ separately: its derivative in $w_\Blk$ is
$\kappa_H\prod_{\Blk'\ne\Blk}(1+v_{\Blk'}+w_{\Blk'})\ge0$, and its derivative in
$v_\Blk$ is at least $\kappa_b\prod_{\Blk'\ne\Blk}(1+v_{\Blk'})\ge0$. Hence, for
chosen upper prices $m_\Blk\ge v_\Blk$ and $n_\Blk\ge w_\Blk$ with
$\sum_\Blk(m_\Blk+n_\Blk)\le U$, it is at most its value at $(m,n)$, which is at
most $\kappa_He^U-(\kappa_H-\kappa_b)e^{\sum_\Blk m_\Blk}$: the expansion of the
latter has the same coefficients $\kappa_b$, $\kappa_H$ on the squarefree
monomials and nonnegative ones on all others. An upper price of the
invariant tags may thus stand in the subtracted exponential, because the form
is monotone in the tag variables; it is never used as a lower bound of a norm
(\cref{lem:oct24:invariant-background,thm:envelopes}).
\end{remark}

\subsection{When a change preserves half-marginals its bound is a variance}

When a comparison law replaces a record by a law with the same actual half
marginals at $\iota$, the difference annihilates the constants on both
sides, and its form against the bracket is a \emph{variance}: the energy of
the record's selected diagonal minus that of the product of its actual half
marginals. This is the interface between the variance credits of Part~V and
the class bounds (CP5): a lower bound for the product energy lowers the
coefficient of the class below $B$.

\begin{lemma}[Variance form at the reserved cut]
\label{lem:record:variance-interface}
Assume (CP1)--(CP3) and (CP8) of \cref{thm:comparison}. Let $X$ be a nonempty
finite set of atoms and, for $\Blk\in X$, let $m_\Blk$ be a probability measure
on the plus-half labels $X_\Blk^{V_+}$ at $\iota$ with
$\operatorname{Diag}_\iota(m_\Blk)$ and $J_\iota(m_\Blk)$ admissible, and
$\Delta_\Blk$ a self-adjoint or skew-adjoint operator on $L^2(m_\Blk)$ with
$\Delta_\Blk\one=\Delta_\Blk^*\one=0$ and $\|\Delta_\Blk\|\le\gamma_\Blk$, the number
of skew-adjoint $\Delta_\Blk$ being even. Put $m_X=\bigotimes_\Blk m_\Blk$,
$\operatorname{Diag}_X=\bigotimes_\Blk\operatorname{Diag}_\iota(m_\Blk)$,
$J_X=\bigotimes_\Blk J_\iota(m_\Blk)$, and let $\zeta$ be the signed measure on
$X_X^V$ whose cut operator at $\iota$ relative to $m_X$ is
$\Delta_X=\bigotimes_\Blk\Delta_\Blk$. Then:
\begin{enumerate}[label=(\roman*)]
\item $Y^{\Nbr}[\operatorname{Diag}_X]\ge Y^{\Nbr}[J_X]$ and
 $|Y_{\zeta\otimes\lambda_{-X}}|\le\prod_\Blk\gamma_\Blk\,
 \bigl(Y^{\Nbr}[\operatorname{Diag}_X]-Y^{\Nbr}[J_X]\bigr)$;
\item $Y^{\Nbr}[\operatorname{Diag}_X]\le(B+o(1))C_XY$ and
 $Y^{\Nbr}[J_X]\ge Y[J_X]-(R_{\rm err}+o(1))C_XY$;
\item consequently, if $Y^{\Nbrz}[J_X]\ge(\ell-o(1))Y$ with $\ell\ge0$ and
 $C_X\le\bar C$, then
 \[
  |Y_{\zeta\otimes\lambda_{-X}}|\le\prod_\Blk\gamma_\Blk\,C_X\bigl(B-g+o(1)\bigr)Y,
  \qquad g=\max\bigl(0,\ell/\bar C-R_{\rm err}\bigr).
 \]
\end{enumerate}
By \ref{it:record:E2}, as in the proof of (ii), $Y^{\Nbrz}[J_X]\ge Y[J_X]$, so
the hypothesis of (iii) holds for every lower bound $\ell$ of $Y[J_X]$; positive
background terms of the expansion of $\Nbrz$ may be added to such a bound
(\cref{prop:background-variance}). Only $\Nbrz\succeq_\iota\mu$ and the priced
difference $\Nbr-\Nbrz$ are used, not an order between $\Nbr$ and $\mu$. The
$o(1)$ consists of energy errors and large-product remainders.
\end{lemma}

\begin{proof}
(i) The number of skew-adjoint factors is even, so $\Delta_X$ is
self-adjoint. As $\Delta_\Blk\one=\Delta_\Blk^*\one=0$, $\Delta_X$ and $\Delta_X^*$
annihilate every vector with a constant factor, in particular the constants
of $L^2(m_X)$; so $\Delta_X=(I-\Jproj)\Delta_X(I-\Jproj)$ and, with
$\gamma=\prod_\Blk\gamma_\Blk\ge\|\Delta_X\|$,
$\gamma(I-\Jproj)\pm\Delta_X\succeq0$. Relative to $m_X$ the cut operators of
$\operatorname{Diag}_X$ and $J_X$ are $I$ and $\Jproj$, so
$\gamma(\operatorname{Diag}_X-J_X)\pm\zeta\succeq_\iota0$, and
\ref{it:record:P5} gives the second assertion; the first is the case
$\Delta_X=0$.

(ii) $\operatorname{Diag}_X$ and $J_X$ are admissible probability measures, so
\ref{it:record:E1} gives the first bound and
$|Y^{\Nbr}[J_X]-Y^{\Nbrz}[J_X]|\le(R_{\rm err}+o(1))C_XY$; moreover
$J_X\succeq_\iota0$, its cut operator being a projection, so
$Y^{\Nbrz}[J_X]\ge Y[J_X]$ by \ref{it:record:E2}.

(iii) By (i), the first bound of (ii), \ref{it:record:E1} for
$\Nbr-\Nbrz$ at $J_X$ and the hypothesis,
$Y^{\Nbr}[J_X]\ge Y^{\Nbrz}[J_X]-(R_{\rm err}+o(1))C_XY\ge(\ell-o(1))Y
-(R_{\rm err}+o(1))C_XY$, and so
$|Y_{\zeta\otimes\lambda_{-X}}|\le\gamma(BC_X+R_{\rm err}C_X-\ell+o(1))Y$. If
$\ell/\bar C\le R_{\rm err}$, use instead $|Y_{\zeta\otimes\lambda_{-X}}|\le\gamma
Y^{\Nbr}[\operatorname{Diag}_X]\le\gamma(B+o(1))C_XY$, valid as
$Y^{\Nbr}[J_X]\ge0$ by \ref{it:record:E2}. Otherwise $\ell\ge C_X\ell/\bar C$,
whence the claim, as $C_X\ge1$.
\end{proof}

\subsection{Energies of the dominators for $\nu$ are at most $Be^TR_I^kY$}

Two places need energies not of $\nu$ itself but of the laws
$M_\Blk$ that dominate it in (CP6): the index-support inequality,
which bounds the signed form $J_\nu$ by such energies with random
inclusion of the index coordinates, and the wrapping terms, which
need the same energies at a fixed interval. The lemma bounds both:
writing $M_\Blk=\Nbr_\Blk+Q_\Blk$, the excess $Q_\Blk$ with total
price $T_Q$ contributes a factor $e^{T_Q}$, and each coupled query
contributes its conductor moment $R_I$.

\begin{lemma}[Envelope energies]\label{lem:envelope-energy}
Assume (CP0), (CP1), (CP2)(b), (CP3) and (CP8) of \cref{thm:comparison}. For
every atom let $M_\Blk=\Nbr_\Blk+Q_\Blk$, with $Q_\Blk$ a nonnegative combination
of admissible positive reserved-diagonal laws and
$\sum_\Blk\|Q_\Blk\|q_\Blk^{2a}\le T_Q$, the prices $\|Q_\Blk\|q_\Blk^{2a}$
satisfying the summability condition of (CP8)(c) (as those of (CP4)(b) and
(CP6) do). Let $\mathcal S\subseteq\mathcal E_k$, and
for $e\in\mathcal S$ let $b_e$ be uniform on $\Z/\Kint\Z$ ($\kappa_e=1$) or
fixed ($\kappa_e=\Kint$) and $s_e$ uniform on $\Z/\Mzero\Z$, all independent.
Let $X$ carry $g^{I_{b_e},\Mzero,s_e}$ at the source $u_e$ of each
$e\in\mathcal S$ and $F_A$ at every other slot of $V_+$, and let
$V=\E_\Jc\E_\omega\ip{X^\Jc}{C_\iota(\bigotimes_\Blk M_\Blk)X^\Jc}$ be its
single-half random-inclusion energy, where $\Jc$ reveals every uniform
interval coordinate and independent prime-power depths with the tails
$\delta_{p,j}$ of (CP0)(d) (for $\mathcal
S=\mathcal E_k$ and uniform intervals, $V=V_{\otimes M}(X)$ of
\cref{def:hyb:coupled}). Then
\[
 V\le\Bigl(Be^{T_Q}\prod_{e\in\mathcal S}\bigl(\kappa_e^{2\theta/d}R_I\bigr)
 +o(1)\Bigr)Y ,
\]
and the same holds with the inputs at the reflected targets $\iota(v_e)$. The
$o(1)$ consists of energy errors and large-product remainders.
\end{lemma}

With $\mathcal S=\mathcal E_k$, uniform intervals and $Q_\Blk=M_{s,\Blk}-\Nbr_\Blk$
of (CP6), this is $V_s\le(Be^TR_I^k+o(1))Y=(2^kJ+o(1))Y$ with
$J=Be^T(R_I/2)^k$, as in \cref{thm:comparison}.

\begin{proof}
Fix $\Jc$ and the intervals. The revealed depths determine $s_e$
modulo $D_{\Jc_e}=\prod_p p^{J_{e,p}}$, a divisor of $\Mzero$. By
\cref{lem:twointerval} and the disjointness of the endpoints, as in
\cref{cor:index-support-frame}(ii), $X^\Jc$ carries
$g^{I_{b_e},D_{\Jc_e},s_e\bmod D_{\Jc_e}}$ at $u_e$ ($e\in\mathcal S$); these support
projections come before any fold, as (O3) requires. By
\eqref{eq:cc:bilinear}, $\ip{X^\Jc}{C_\iota(\bigotimes_\Blk M_\Blk)X^\Jc}$ is the
energy of $\bigotimes_\Blk M_\Blk$ with the input of $u$ at $u$ and at $\iota u$
($u\in V_+$): each $e\in\mathcal S$ has the two changed slots $u_e$, $\iota u_e$,
each of coarse price $(\Kint\sq(D_{\Jc_e}))^{a/d}\kappa_e^{\theta/d}$
(\cref{def:record:structured-energy}, with the square step of
\cref{lem:record:fibre-scale}).
Expand the outside laws fully and split by the total selected union
as above; the large unions contribute one aggregate $o(1)Y$.
For the small unions, start with
$\bigotimes_\Blk M_\Blk=\sum_Z\bigotimes_{\Blk\in Z}Q_\Blk\otimes\Nbr_{-Z}$ over the
finite sets $Z$ of atoms. Each $\bigotimes_{\Blk\in Z}Q_\Blk$ is admissible and
positive, so the small-union part of the proof of \cref{lem:bracket-energy}
bounds its contribution by
$(B+o(1))\prod_{\Blk\in Z}\|Q_\Blk\|\,\Pi_ZY$, where
$\Pi_Z=C_Z\prod_{e\in\mathcal S}(\Kint\sq(D_{\Jc_e}))^{2a/d}\kappa_e^{2\theta/d}$.
Since $\sum_Z\prod_{\Blk\in Z}\|Q_\Blk\|q_\Blk^{2a}=\prod_\Blk(1+\|Q_\Blk\|q_\Blk^{2a})
\le e^{T_Q}$, the energy at fixed $\Jc$ is at most
$(B+o(1))e^{T_Q}\prod_{e\in\mathcal S}(\Kint\sq(D_{\Jc_e}))^{2a/d}
\kappa_e^{2\theta/d}\,Y$. In \cref{lem:bracket-energy} the conductors stay
fixed through every fold and geometric mean, and complementary folds keep
two occurrences of each conductor (\cref{lem:record:folded-energy}(e), (O4)).
Only now average over $\Jc$: the $\Jc_e$ are independent with the law of
(CP0)(d), so $\E_\Jc\prod_{e\in\mathcal S}(\Kint\sq(D_{\Jc_e}))^{2a/d}
=R_I^{|\mathcal S|}$, and the errors are uniform in $\Jc$. For the reflected
targets use the slots $\iota(v_e)$ and $v_e$.
\end{proof}

\begin{remark}[The trivial mean credit and the bookkeeping of the closing]
\label{lem:record:domination}
Assume (CP1)--(CP3), (CP6) and (CP8). For every atom and
$s\in\{{\rm L},{\rm R}\}$,
$M_{s,\Blk}\pm\lambda_\Blk=(\Nbr_\Blk\pm\lambda_\Blk)+(M_{s,\Blk}-\Nbr_\Blk)
\succeq_\iota0$: the first summand by (CP3) (off the first-stage atoms
$\Nbr_\Blk\pm\lambda_\Blk$ is $2\mu_\Blk$ or $0$), the second by
\ref{it:record:E3}. By \cref{lem:record:block-gram}(d),
$\bigotimes_\Blk M_{s,\Blk}\pm\bigotimes_\Blk\lambda_\Blk\succeq_\iota0$, and at
$\bar X$ this is $e_s\ge Y_\lambda\ge(1-L_0-o(1))Y$
(\cref{lem:record:first-stage}). Hence (CP7) holds with the trivial credit
$m=1-L_0$, and then
$\Theta=(1-2^{-k})(L_0+\mathcal E)+J+W+2^{-k}L_0=L_0+(1-2^{-k})\mathcal E+J+W$:
the first-stage loss enters $2^k-1$ times through $Y_\nu$ and once through
the mean credit, so it keeps coefficient one, as do the matching term and
the wrap, and only the reserved cost drops to $1-2^{-k}$. No resource is
counted twice: the class bounds and credits enter only through the reserved
cost, a separate inequality between $Y_\nu$ and $Y_\lambda$
(\cref{lem:record:record-expansion}); the matching prices enter only through
the bounds for $V_{\rm L}$, $V_{\rm R}$ (\cref{lem:envelope-energy}); and this
credit uses no price at all, only positivity, and no order between $\Nbr$ and
$\mu$. The construction uses the larger credit of
\cref{prop:record:mean-credit} for (CP7); this remark is not one of its
premises.
\end{remark}

\paragraph{Consequences and what to read next.}
From Part~II: cut matrices, their order and the Kronecker rule
(\cref{def:cc:order}, \eqref{eq:cc:bilinear}, \eqref{eq:cc:kron},
\cref{def:rp}, \cref{rem:cut-symmetries}), the common-cut order (\cref{lem:commoncut}), the block-Gram
calculus (\cref{lem:record:block-gram}) and the sector lemma
(\cref{lem:sector}); the functional, its inputs, the index averages and the
reference hypotheses (\cref{def:record:functional}, \cref{lem:twointerval},
\cref{hyp:record:reference}); structured energies and their estimate, folded
energies, admissibility, the square step, the cancellation and the large
products
(\cref{def:record:structured-energy,prop:record:density-structured,%
lem:record:folded-energy,def:admissible,lem:record:fibre-scale,%
lem:record:density-can,lem:record:density-lp}), at a least normalized
counterexample (\cref{def:record:counterexample}). From Part~III: the frame
data and hypotheses of \cref{thm:comparison}, the index law, the nonvanishing square classes
and the signed index form (\cref{def:hyb:coupled}, \cref{lem:func:unitsq},
\eqref{eq:record:expansion}), the
comparison law (\cref{def:record:comparison-law}) and the prices
(\cref{def:price}).
